\subsection{Stone Pile за $\mathcal{O}(N^2\cdot W)$}

\subsubsection{Описание решения}

Воспользуемся динамическим программированием для решения задачи, а именно приравняем её к задаче о рюкзаке. Введём следующие обозначения:
\begin{itemize}
  \item \( f(w,i) \) --- минимальное число камней среди первых \( i \) камней в массиве, у которых масса в сумме хотя бы \( w \).
  \item \( W \) --- максимально возможная масса камня.
  \item \( W_{\text{all}} \) --- сумма масс камней данного массива, \( W_{\text{all}}\leq N\cdot W \).
\end{itemize}

Для базы рекуррентной формулы положим:
\[
  f(w,0)=
  \begin{cases} 
    0, & \text{если } w = 0 \\
    N + 1, & \text{если } w > 0
  \end{cases}
\]

Тогда шаг рекурренты выглядит следующим образом:
\[
  f(w, i) = 
\begin{cases}
\min\bigl(f(w, i-1),\; f(w-w_i, i-1) + 1\bigr), & \text{если } w \geq w_i \\[4pt]
f(w, i-1), & \text{если } w < w_i
\end{cases}
\]

После вычисления значений функции, ответ будет во множестве:
\[
  X=\left\{f(w,N)\mid \forall w=\overline{0,W_{\text{all}}},\ f(w,N)\neq N+1\right\}
\]

Ответ рассчитывается по формуле:
\[
  x=\min\left\{\abs{2w-W_{\text{all}}}\mid f(w,N)\in X\right\}
\]

Напоследок можно заметить, что рекуррентное выражение использует для подсчёта лишь предыдущий слой относительно индекса \( i \). С этим наблюдением можно оптимизировать использование памяти и хранить вместо двумерного массива только два слоя: текущий и предыдущий.

\subsubsection{Асимптотическая сложность}

\textit{Временная сложность решения} --- \( \mathcal{O}(N^2\cdot W) \). Для расчёта ответа высчитывается \( N\cdot W_{\text{all}}\leq N^2\cdot W \) значений рекуррентного выражения. Ответ определяется на последнем шаге за линейное время.

\textit{Пространственная сложность решения} --- \( \mathcal{O}(N+W) \). Для вычисления значений рекуррентного выражения используется два массива длины \( W_{\text{all}}\leq W \), исходные данные хранятся в массиве длины \( N \).

\subsubsection{Листинг кода}

\begin{lstlisting}
#include <iostream>
#include <vector>

int main() {
  int N;
  std::cin >> N;
  std::vector<int> arr;
  long total_sum = 0;
  for (int i = 0; i < N; i++) {
    int item;
    std::cin >> item;
    arr.push_back(item);
    total_sum += item;
  }
  std::vector<int> min_coins[2];
  int INF = N + 1;
  for (int j = 0; j < 2; j++) {
    std::vector<int> step;
    for (int i = 0; i <= total_sum; i++) {
      step.push_back(INF);
    }
    step[0] = 0;  // dp base
    min_coins[j] = step;
  }
  short curr_step_idx = 0, prev_step_idx = 1;
  for (int i = 1; i <= N; i++) {
    curr_step_idx ^= 1;
    prev_step_idx ^= 1;
    for (long cost = 0; cost <= total_sum; cost++) {
      int prev = min_coins[prev_step_idx][cost];
      long prev_cost = cost - arr[i - 1];
      int curr = prev_cost >= 0 ? min_coins[prev_step_idx][cost - arr[i - 1]] + 1 : INF;
      min_coins[curr_step_idx][cost] = prev <= curr ? prev : curr;
    }
  }
  long diff = total_sum;
  for (long cost = 0; cost < total_sum; cost++) {
    long curr_diff = abs(2 * cost - total_sum);
    if (min_coins[curr_step_idx][cost] != INF && curr_diff < diff) {
      diff = curr_diff;
    }
  }
  std::cout << diff;
  return 0;
}
\end{lstlisting}